\chapter{Annexes}

\section{Commandes d'installation et d'exécution}
Les commandes suivantes permettent de mettre en place l'environnement de développement :

\begin{lstlisting}[language=bash,caption={Installation et exécution de l'application}]
# Creation de l'environnement virtuel
python -m venv venv
source venv/bin/activate  # Linux/Mac
venv\Scripts\activate     # Windows

# Installation des dependances
pip install -r requirements.txt

# Application des migrations
python manage.py migrate

# Creation des donnees de demonstration
python manage.py seed_data

# Lancement du serveur de developpement
python manage.py runserver
\end{lstlisting}

\section{Identifiants de démonstration}
Après exécution de la commande \texttt{seed\_data}, les comptes suivants sont disponibles :

\begin{table}[H]
\centering
\caption{Comptes de démonstration}
\begin{tabularx}{\textwidth}{|l|l|X|}
\hline
\textbf{Profil} & \textbf{Identifiant} & \textbf{Mot de passe} \\
\hline
Administrateur & admin & admin1234 \\
\hline
Employé & employe & employe1234 \\
\hline
Client & client & client1234 \\
\hline
\end{tabularx}
\end{table}

\section{Extrait des routes principales}
\begin{lstlisting}[language=Python,caption={Routes principales de l'application}]
path("", views.accueil, name="accueil")
path("login/", views.login_view, name="login")
path("logout/", views.logout_view, name="logout")
path("dashboard/", views.dashboard, name="dashboard")
path("password-reset/", views.password_reset_view, name="password_reset")
path("voitures/", views.liste_voitures, name="liste_voitures")
path("voitures/ajouter/", views.ajouter_voiture, name="ajouter_voiture")
path("clients/", views.liste_clients, name="liste_clients")
path("reservations/", views.liste_reservations, name="liste_reservations")
path("client/dashboard/", views.client_dashboard, name="client_dashboard")
path("client/voitures/", views.client_voitures, name="client_voitures")
path("client/reserver/", views.client_creer_reservation, name="client_creer_reservation")
path("roles/", views.gestion_roles, name="gestion_roles")
path("statistiques/", views.statistiques, name="statistiques")
path("historique/", views.historique, name="historique")
path("notifications/", views.liste_notifications, name="liste_notifications")
path("rapport/pdf/", views.rapport_pdf, name="rapport_pdf")
\end{lstlisting}

\section{Extrait de calcul métier}
\begin{lstlisting}[language=Python,caption={Calcul automatique du montant total}]
def calculer_montant(self):
    nombre_jours = (self.date_fin - self.date_debut).days
    if nombre_jours <= 0:
        nombre_jours = 1
    return nombre_jours * self.voiture.prix_jour
\end{lstlisting}

\section{Exemple de contrôle de conflit}
\begin{lstlisting}[language=Python,caption={Detection des chevauchements de reservation}]
conflit = Reservation.objects.filter(
    voiture=reservation.voiture,
    statut__in=["en_attente", "en_cours"],
    date_debut__lt=reservation.date_fin,
    date_fin__gt=reservation.date_debut,
)
\end{lstlisting}

\section{Extrait de journalisation}
\begin{lstlisting}[language=Python,caption={Fonction d'enregistrement dans l'historique}]
def enregistrer_historique(action, description, 
                           utilisateur=None, reservation=None):
    Historique.objects.create(
        action=action,
        description=description,
        utilisateur=utilisateur,
        reservation=reservation,
    )
\end{lstlisting}

\section{Structure du fichier requirements.txt}
\begin{lstlisting}[language=bash,caption={Dependances du projet}]
Django>=6.0,<6.1
reportlab>=4.0
Pillow>=11.0
\end{lstlisting}

\section{Variables d'environnement}
Le fichier \texttt{.env.example} fournit le modèle de configuration suivant :

\begin{lstlisting}[language=bash,caption={Modele de fichier .env.example}]
SECRET_KEY=votre-cle-secrete-ici
DEBUG=True
ALLOWED_HOSTS=localhost,127.0.0.1
\end{lstlisting}

\clearpage
